\noindent\textbf{(a)} \(k=27\pi^2/5\,\mathrm{N/m}\), \(c=2\pi\sqrt{14}/5\,\mathrm{N\,s/m}\).

\noindent\textbf{(b)} \[m_\pm=\frac{k\pm\sqrt{k^2-c^2(2\pi f)^2}}{2(2\pi f)^2}.\]
Hay dos masas si \(0<f<k/(2\pi c)\), una si hay igualdad y ninguna por encima del límite.

\noindent\textbf{(c)} \(m=(135+\sqrt{18001})/8\,\mathrm{kg}\simeq33.646\,\mathrm{kg}\).
